\chapter{Theoretische und technologische Grundlagen}

\section{Künstliche Intelligenz und Machine Learning}

Dieses Kapitel vermittelt die technologischen Grundlagen, die für das Verständnis von KI-gestütztem Monitoring und algorithmischer Governance notwendig sind.

\subsection{Grundbegriffe}

Künstliche Intelligenz (KI) befasst sich im Kern mit der Frage, wie Maschinen eine Welt, die weit größer und komplizierter ist als sie selbst, wahrnehmen, verstehen, vorhersagen und manipulieren können.\vglfootcite[1\,ff.]{RussellNorvig2021} Ein zentrales Konzept ist hierbei der rationale Agent, der seine Umwelt durch Sensoren wahrnimmt und über Aktuatoren darauf einwirkt. Ein solcher Agent gilt als rational, wenn er für jede mögliche Wahrnehmungssequenz eine Aktion wählt, die unter Berücksichtigung seines Wissens das erwartete Leistungsmaß maximiert.\vglfootcite[34]{RussellNorvig2021}

Machine Learning (ML) ist ein Teilbereich der KI, der es Systemen ermöglicht, sich an neue Umstände anzupassen sowie Muster in Daten zu erkennen und zu extrapolieren.\vglfootcite[105]{Goodfellow2016} Die Quellen unterscheiden dabei drei wesentliche Lernparadigmen:

\begin{itemize}
    \item \textbf{Überwachtes Lernen (Supervised Learning):} Das System lernt Funktionen auf Basis von markierten Beispieldaten (Inputs und Zielvorgaben).\vglfootcite[105]{Goodfellow2016}
    \item \textbf{Unüberwachtes Lernen (Unsupervised Learning):} Hierbei werden Kategorien oder Strukturen in Sammlungen nicht etikettierter Daten erkannt, beispielsweise durch Clustering.\vglfootcite[105]{Goodfellow2016}
    \item \textbf{Bestärkendes Lernen (Reinforcement Learning):} Ein Agent lernt durch Interaktion mit einer Umgebung, indem er für seine Aktionen Belohnungen oder Bestrafungen erhält, um eine optimale Strategie zu entwickeln.\vglfootcite[1\,ff.]{SuttonBarto2018}
\end{itemize}

\subsection{Überwachungs- und Analyseverfahren}

Monitoring in digitalen Ökosystemen basiert auf der Interpretation von Sensordaten, die Aspekte der Umgebung in eine für Computer verwertbare Form übersetzen.\footnote{TODO: Quelle einfügen} Ein wesentliches Verfahren ist die Datenassoziation, bei der mehrere Beobachtungen kombiniert werden, um Objekte in einer komplexen Welt konsistent zu verfolgen.

Die erhobenen Daten werden häufig zu digitalen Nutzerprofilen zusammengeführt. Diese Profile enthalten Informationen über Interessen, Verhaltensmuster, Präferenzen sowie Wahrscheinlichkeiten zukünftiger Handlungen. Nutzerprofile bilden die Grundlage vieler datengetriebener Analyse-, Prognose- und Empfehlungssysteme, da sie eine individualisierte Bewertung und Vorhersage von Nutzerverhalten ermöglichen.\footnote{TODO: Quelle einfügen}

Für die Analyse werden häufig Hidden-Markov-Modelle (HMM) eingesetzt, die sich besonders für zeitliche oder dynamische Prozesse eignen, bei denen die eigentlichen Zustände nicht direkt beobachtbar (\enquote{hidden}) sind, sondern nur über stochastisch abhängige Beobachtungen erschlossen werden können.\footnote{TODO: Quelle einfügen} Im Bereich des unüberwachten Lernens spielt die Ausreißererkennung (Outlier Detection) eine wichtige Rolle, um Datenpunkte zu identifizieren, die durch Rauschen, Messfehler oder ungewöhnliches Verhalten stark von der Norm abweichen.\footnote{TODO: Quelle einfügen}

Ein weiteres zentrales Analyseverfahren stellen Scoring- und Ranking-Systeme dar. Hierbei werden Personen, Inhalte oder Objekte anhand definierter Merkmale bewertet und in eine Rangfolge eingeordnet. Die daraus resultierenden Bewertungen dienen häufig als Grundlage für automatisierte Entscheidungen, Personalisierungsmechanismen und die Priorisierung von Informationen innerhalb digitaler Plattformen.\footnote{TODO: Quelle einfügen}

\subsection{Empfehlungssysteme, Predictive Analytics und Verhaltenssteuerung}

Predictive Analytics nutzt maschinelle Intelligenz, um Vorhersagen über menschliches Verhalten zu erstellen, wobei diese Vorhersagen umso wertvoller werden, je näher sie der Gewissheit kommen.\footnote{TODO: Quelle einfügen} Diese Systeme basieren oft auf Wahrscheinlichkeitsmodellen wie Bayesschen Netzen, die Abhängigkeiten zwischen Variablen darstellen, um auch unter Unsicherheit fundierte Schlüsse ziehen zu können.\footnote{TODO: Quelle einfügen}

Die Grundlage moderner Predictive-Analytics-Systeme besteht häufig in der kontinuierlichen Auswertung großer Mengen historischer und aktueller Verhaltensdaten. Ziel ist die Identifikation von Mustern, Zusammenhängen und Wahrscheinlichkeiten zukünftiger Handlungen. Die Qualität solcher Vorhersagen steigt dabei in der Regel mit der Menge und Aktualität der verfügbaren Daten.\footnote{TODO: Quelle einfügen}

Empfehlungssysteme sind eine spezifische Anwendung, bei der Algorithmen die Präferenzen eines Nutzers mit denen anderer abgleichen, um relevante Informationen oder kulturelle Inhalte vorzuschlagen.\footnote{TODO: Quelle einfügen}

Empfehlungssysteme basieren häufig auf Verfahren des kollaborativen Filterns, der Inhaltsanalyse oder hybriden Ansätzen. Dabei werden Nutzerpräferenzen analysiert, um relevante Inhalte, Produkte oder Informationen vorzuschlagen. Solche Systeme finden Anwendung in sozialen Netzwerken, Streaming-Plattformen, Suchmaschinen und E-Commerce-Systemen.\footnote{TODO: Quelle einfügen}

Zwischen Nutzer und Plattform entstehen dabei kontinuierliche Rückkopplungsschleifen. Jede Interaktion eines Nutzers erzeugt neue Daten, welche zur Verbesserung zukünftiger Vorhersagen und Empfehlungen genutzt werden. Dieser Prozess ermöglicht eine fortlaufende Optimierung algorithmischer Modelle und eine zunehmend präzisere Analyse individuellen Verhaltens.\footnote{TODO: Quelle einfügen}

In der Logik des sogenannten Überwachungskapitalismus werden solche Systeme genutzt, um aus \enquote{Verhaltensüberschüssen} Vorhersageprodukte zu fertigen, die in speziellen Märkten für zukünftiges Verhalten gehandelt werden.\footnote{TODO: Quelle einfügen} Das Ziel ist hierbei oft nicht nur die Vorhersage, sondern die Verhaltensmodifikation, um garantierte kommerzielle Ergebnisse zu erzielen.\footnote{TODO: Quelle einfügen}

Aus Sicht digitaler Plattformen beschränkt sich die Nutzung solcher Systeme nicht ausschließlich auf die Vorhersage zukünftigen Verhaltens. Vielmehr können algorithmische Empfehlungen, Rankings und Personalisierungsmechanismen genutzt werden, um Aufmerksamkeit zu lenken, Interaktionen zu fördern und bestimmte Verhaltensweisen wahrscheinlicher zu machen. Die Übergänge zwischen Verhaltensprognose und Verhaltensbeeinflussung bilden dabei einen zentralen Untersuchungsgegenstand dieser Arbeit.\footnote{TODO: Quelle einfügen}

\section{Digitale Informationssysteme}

Das Verständnis moderner Monitoring-Systeme setzt eine fundierte Auseinandersetzung mit der theoretischen Basis digitaler Informationssysteme voraus. In der Wirtschaftsinformatik werden diese Systeme nicht rein technologisch betrachtet, sondern als komplexe Gefüge, die weit über Hardware- und Softwarekomponenten hinausgehen. Die folgenden Abschnitte erläutern die grundlegenden Systemarchitekturen, die internen Datenflüsse sowie die wachsende Bedeutung von Plattformökosystemen als Rahmen für datengetriebene Dienste.

\subsection{Systemarchitekturen}

Ein digitales Informationssystem lässt sich grundlegend als ein Satz miteinander verbundener Komponenten definieren, die Daten sammeln, verarbeiten, speichern und verteilen, um Entscheidungsprozesse und Kontrollmechanismen in einer Organisation zu unterstützen.\footnote{TODO: Quelle einfügen} Der theoretische Kern dieser Definition liegt in der Betrachtung als soziotechnisches System. In diesem Ansatz wird das Informationssystem nicht isoliert als Technologie verstanden, sondern als ein Zusammenspiel der Dimensionen Organisation, Management und Technologie. Die Leistung eines solchen Systems wird dann optimiert, wenn sich Technologie und Organisation gegenseitig anpassen, bis eine zufriedenstellende Übereinstimmung erreicht ist. Dabei bilden Menschen, ihre Rollen und die politische Kultur innerhalb einer Organisation die soziale Komponente, während Daten und Infrastruktur die technologische Basis darstellen.\footnote{TODO: Quelle einfügen}

Die technische Realisierung dieser Systeme erfolgt über verschiedene Architekturmuster. Historisch bedeutsam ist die Client-Server-Architektur, bei der spezialisierte Rechnerplattformen (Server) Funktionen wie Datenbankmanagement oder Programmausführung für andere Computer im Netzwerk (Clients) bereitstellen. Diese Architektur minimiert den Netzwerkverkehr, da nur die für eine spezifische Abfrage benötigten Daten übertragen werden, und ermöglicht eine zentrale Implementierung von Sicherheitsmechanismen auf dem Server. Moderne Erweiterungen führen zu Mehrschichtarchitekturen (Multi-Tier Architecture), in denen beispielsweise Applikationsserver zwischen den browserbasierten Clients und den Backend-Datenbanken agieren, um die Lastverteilung und Wartbarkeit zu erhöhen.\footnote{TODO: Quelle einfügen}

Einen weiteren Evolutionsschritt stellt die serviceorientierte Architektur (SOA) dar. Sie nutzt modulare Anwendungsdienste, die es Systemen ermöglichen, standardisiert miteinander zu interagieren. Diese Modularität erlaubt es Unternehmen, bestehende Anwendungen als Bausteine für neue, funktionsübergreifende Prozesse zu nutzen, was die Flexibilität und Skalierbarkeit massiv erhöht. In einer noch feingranulareren Form finden sich Microservices, die durch eine erweiterbare Codebasis Kernfunktionalitäten bereitstellen, welche von spezialisierten Zusatzmodulen (Apps) über definierte Schnittstellen (APIs) angesprochen werden.\footnote{TODO: Quelle einfügen}

In der aktuellen Praxis dominieren zunehmend Cloud-basierte Informationssysteme. Hierbei werden IT-Ressourcen wie Rechenleistung, Speicher und Anwendungen über das Internet als Abonnementservice bereitgestellt. Cloud-Modelle bieten eine hohe Skalierbarkeit, also die Fähigkeit, die Verarbeitungskapazität dynamisch an steigende Nutzerzahlen oder Transaktionsvolumina anzupassen, ohne massiv in physische Infrastruktur investieren zu müssen. Diese modernen Architekturen bilden das technologische Rückgrat für die Verarbeitung jener massiven Datenmengen, die in Monitoring-Szenarien anfallen.\footnote{TODO: Quelle einfügen}

Die strukturelle Beschaffenheit dieser Architekturen bestimmt maßgeblich, wie Informationen innerhalb des Systems fließen und welche Logiken zur Verarbeitung dieser Daten herangezogen werden können.

\subsection{Datenflüsse und Entscheidungslogiken}

Innerhalb digitaler Informationssysteme durchlaufen Daten einen definierten Lebenszyklus, der mit der Datenerfassung (Input) beginnt. Dieser Prozess umfasst das Sammeln und Erfassen von Rohdaten aus der internen oder externen Umgebung des Systems. Eine effiziente Methode ist hierbei die Quellendatenautomatisierung (Source Data Automation), bei der Daten unmittelbar am Ort ihrer Entstehung in digitaler Form erfasst werden, beispielsweise durch Sensoren oder mobile Endgeräte. Die anschließende Datenverarbeitung (Processing) transformiert diese Rohdaten in Informationen, indem sie klassifiziert, sortiert, berechnet oder formatiert werden, um ihnen eine Bedeutung für den menschlichen Nutzer zu verleihen.\footnote{TODO: Quelle einfügen}

Die Gewinnung von Informationen dient primär der Entscheidungsunterstützung. Während operative Systeme (Transaction Processing Systems -- TPS) tägliche Routinetransaktionen verarbeiten, aggregieren Management-Informationssysteme (MIS) diese Daten zu Berichten, um die Effizienz der Betriebsabläufe zu überwachen. Für komplexere, unstrukturierte Problemstellungen werden Decision Support Systems (DSS) eingesetzt. Ein DSS kombiniert Daten mit anspruchsvollen Analysemodellen, um \enquote{Was-wäre-wenn}-Szenarien durchzuspielen und Managern bei individuellen, nicht-routinemäßigen Entscheidungen zu assistieren.\footnote{TODO: Quelle einfügen}

Im Bereich der Business Intelligence (BI) werden diese Konzepte weiter gefasst. BI umfasst Prozesse und Softwaretools zur Sammlung, Analyse und Bereitstellung von Daten, um fundierte organisatorische Entscheidungen zu treffen. Durch Techniken wie Predictive Analytics können auf Basis historischer Daten und Annahmen über zukünftige Bedingungen Vorhersagen über Ereignisse, wie etwa künftige Umsätze oder das Ausfallrisiko von Kunden, getroffen werden. Ein wesentlicher Trend ist hierbei die zunehmende Automatisierung. Automatisierte Entscheidungsprozesse nutzen Algorithmen und mathematische Modelle, um in Echtzeit auf Probleme oder Chancen zu reagieren (Sense and Respond).\footnote{TODO: Quelle einfügen}

Für KI-gestützte Monitoring-Systeme sind diese Entscheidungslogiken von zentraler Bedeutung. Da menschliche Akteure in hochvolumigen Datenumgebungen oft als Hindernis für die Effizienz angesehen werden, sind Systeme so konzipiert, dass sie weitgehend autonom agieren. Die Qualität der Entscheidungsunterstützung hängt dabei direkt von der Fähigkeit des Systems ab, Muster in den Datenströmen zu erkennen und diese in präskriptive Handlungsanweisungen zu übersetzen. Der Zusammenhang zwischen systematischer Datenanalyse und organisatorischem Erfolg wird besonders deutlich, wenn Organisationen \enquote{auf Basis von Analysen konkurrieren} (Competing on Analytics), indem sie jeden Aspekt ihrer Geschäftsprozesse optimieren.\footnote{TODO: Quelle einfügen}

Diese internen Logiken der Datenverarbeitung und Entscheidung finden ihren stärksten Ausdruck heute in vernetzten Infrastrukturen, die über einzelne Organisationen hinausgehen.

\subsection{Plattformökosysteme}

Die moderne Informationsökonomie ist geprägt durch den Aufstieg von digitalen Plattformen, die als Geschäftsmodelle fungieren, welche wertschöpfende Interaktionen zwischen externen Produzenten und Konsumenten ermöglichen. Plattformen stellen eine offene Infrastruktur bereit und definieren die Regeln für diesen Austausch. In der Literatur wird zwischen zwei Grundtypen unterschieden: Innovationsplattformen, die technologische Bausteine für Drittanbieter bereitstellen (z.\,B. iOS, Android, AWS), und Transaktionsplattformen, die als Vermittler auf Marktplätzen agieren (z.\,B. Amazon Marketplace, Facebook, Uber).\footnote{TODO: Quelle einfügen}

Ein entscheidendes Merkmal von Plattformen ist das Auftreten von Netzwerkeffekten. Diese positiven Rückkopplungsschleifen führen dazu, dass der Wert einer Plattform für jeden einzelnen Nutzer steigt, je mehr andere Teilnehmer dem Netzwerk beitreten. Dabei wird zwischen direkten Effekten (Nutzer ziehen Nutzer an) und indirekten bzw. kreuzseitigen Effekten (eine Marktseite zieht die andere an, z.\,B. App-Entwickler und Smartphone-Nutzer) differenziert. Diese Dynamik verleiht Plattformen eine natürliche Tendenz zur Monopolisierung und ermöglicht ein exponentielles Wachstum, da sie \enquote{Asset-light} operieren -- also Werte schöpfen, ohne die zugrunde liegenden Ressourcen (wie Fahrzeuge bei Uber oder Immobilien bei Airbnb) selbst besitzen zu müssen.\footnote{TODO: Quelle einfügen}

In diesem Kontext werden Daten zur zentralen Ressource. Plattformen sind als \enquote{extraktiver Apparat} konzipiert, der kontinuierlich Daten aus Nutzeraktivitäten, natürlichen Prozessen und Transaktionen gewinnt. Diese Daten dienen nicht nur der Personalisierung von Diensten, sondern verbessern durch maschinelles Lernen auch stetig die Algorithmen für Suche und Empfehlungen, was wiederum die Plattform attraktiver macht. Die Plattform-Governance spielt hierbei eine Schlüsselrolle, da der Plattformbetreiber die Regeln festlegt, nach denen Produkte entwickelt und Interaktionen bewertet werden.\footnote{TODO: Quelle einfügen}

Beispiele wie Google, YouTube, Facebook, TikTok und Amazon illustrieren, wie Plattformökosysteme als Grundlage datengetriebener Dienste fungieren. Sie nutzen Algorithmen, um die Aufmerksamkeit der Nutzer zu steuern, Inhalte zu ranken und das Verhalten durch gezielte Empfehlungen zu beeinflussen. Die massenhafte Sammlung digitaler Fußabdrücke ermöglicht die Erstellung umfassender Nutzungsprofile, die zur Grundlage für algorithmische Entscheidungen über die Sichtbarkeit von Informationen werden.\footnote{TODO: Quelle einfügen}

Die in diesem Kapitel dargelegten Grundlagen der Systemarchitekturen, Datenflüsse und Plattformlogiken verdeutlichen, dass digitale Informationssysteme die infrastrukturelle Voraussetzung für modernes Monitoring schaffen. Der Übergang von der rein funktionalen Unterstützung organisatorischer Prozesse hin zu einer umfassenden, algorithmisch gesteuerten Verhaltenssteuerung führt direkt zum nächsten Kapitel, welches sich spezifisch mit den Mechanismen von Monitoring- und Bewertungssystemen befasst.

\section{Monitoring- und Bewertungssysteme}
\subsection{Datenerfassung}
\subsection{Analyse und Modellbildung}
\subsection{Algorithmische Bewertung und Klassifikation}

\section{Methodisches Vorgehen}

\subsection{Forschungsansatz}
This thesis follows a literature-based research approach with a qualitative and interdisciplinary orientation. The objective of the thesis is to examine how AI-supported monitoring systems function as digital information systems and to analyze how these systems enable data-driven behavioral analysis and influence within digital environments.

Due to the socio-technical nature of the topic, the research combines perspectives from business informatics, information systems research, artificial intelligence, digital platform economics, and social sciences. The focus is placed on the technological structures, data processing mechanisms, and socio-economic implications of AI-supported monitoring and recommendation systems.

The thesis does not aim to empirically measure user behavior through experiments or surveys. Instead, existing scientific literature, technical publications, and selected practical examples are analyzed and critically compared.

\subsection{Literaturrecherche}
The literature research process was conducted iteratively and thematically. At the beginning of the research process, exploratory searches were performed in order to identify relevant concepts, technological mechanisms, and central authors within the research domains of artificial intelligence, digital monitoring systems, recommendation algorithms, platform ecosystems, and behavioral influence.

The identification of relevant literature was supported by AI-assisted exploratory search methods as well as academic search platforms and scientific databases. During the research process, scientific publications were continuously expanded through citation tracking and thematic cross-referencing. In particular, references and citations within already identified scientific papers were used to locate additional relevant literature and foundational works.

The literature selection focused primarily on peer-reviewed journal articles, conference papers, scientific books, and academically reputable publications within the fields of Information Systems, Artificial Intelligence, platform economics, and digital governance.

\subsection{Auswahlkriterien der Quellen}
The selection of literature was based on several quality criteria. Preference was given to scientific and peer-reviewed publications with thematic relevance to AI-supported monitoring systems, digital information systems, algorithmic decision-making, recommendation systems, surveillance mechanisms, and platform-based behavioral analysis.

In addition to scientific literature, selected reports, documentaries, and journalistic sources were included in limited form where they contributed practical examples or illustrated current societal debates. Such sources were not treated as scientific evidence but rather as supplementary contextual material.

Particular attention was paid to the scientific reputation of authors, publication quality, citation frequency, topical relevance, and methodological transparency of the respective sources.

\subsection{Methodische Grenzen}
The research topic contains several methodological limitations. Many modern AI-supported monitoring and recommendation systems operate as proprietary systems, meaning that internal algorithms, data structures, and optimization mechanisms are often not fully transparent to external researchers.

Furthermore, discussions surrounding digital surveillance, platform governance, and state-supported monitoring systems are frequently influenced by political, ideological, and media-related perspectives. As a result, the available literature and public information may contain normative interpretations and differing viewpoints.
